\section{Experimental methodology}
\label{sec:setup}

\subsection{Platform and software}

\Cref{tab:platform} lists the platform. ``Default'' means the stock
governors: \texttt{nvhost\_podgov} for the GPU (eight
steps from 306 to 1020\,MHz) and \texttt{schedutil} for the CPU.
``Locked'' means \texttt{jetson\_clocks}, which pins the CPU, GPU and
memory controller (EMC) clocks to their
maxima~\cite{nvidia2024jetsonperf}. Readback at ${\approx}4$\,Hz, idle
and under a 15\,s inference load, confirmed the EMC pin (constant
3199\,MHz). Each locked phase
saves the default configuration to a fresh file and restores it
afterwards; governor and clock limits were read back after every phase. The two NumPy versions in \cref{tab:platform} differ
because Ultralytics 8.4.172 requires the newer release. This affects
only the short COCO row-conversion step, whose code is identical in
the custom pipelines and the baseline.

\begin{table}[t]
\centering
\caption{Platform and software. BSP is the board support package.}
\label{tab:platform}
\footnotesize
\begin{tabular}{@{}ll@{}}
\toprule
Board & Jetson Orin Nano Super Developer Kit, \\ & 8\,GB module \\
CPU & 6$\times$ Cortex-A78AE, two frequency \\
    & domains (CPUs 0--3, 4--5), 115--1728\,MHz \\
GPU & Ampere, 1024 cores, 306--1020\,MHz \\
Power mode & \texttt{MAXN\_SUPER} \\
OS / BSP & Ubuntu 22.04, Jetson Linux R36.5.2 \\
Kernel & 5.15.199-tegra \\
Cooling & fan under thermal control \\
 & (\texttt{nvfancontrol}); locked runs pin CPU, \\
 & GPU and EMC clocks, not the fan \\
CUDA / TensorRT & 12.6 / 10.3.0 \\
Detector software & Ultralytics 8.4.172 \\
Baseline runtime & PyTorch 2.5.0 (NVIDIA JetPack build) \\
Host libraries & Python 3.10.12, OpenCV 4.10.0, \\ & NumPy 1.21.5 (custom), 1.26.4 (Ultralytics) \\
Power sensor & INA3221 \texttt{VDD\_IN} (module input), \\ & 20\,Hz \\
\bottomrule
\end{tabular}
\end{table}

\subsection{Replication platforms}
\label{sec:replplatforms}

Two further boards (\cref{sec:crossplatform}) rerun
the core protocol with the smallest changes their stacks require. On
the Raspberry Pi~5, the model is the Hailo Model Zoo's precompiled
HEF (Hailo Executable Format) \texttt{yolov8n\_h8.hef}~\cite{hailo2023modelzoo} for the
Hailo-8 (INT8, NMS on chip, score threshold 0.2 fixed at compile time),
run through HailoRT~4.23.0 asynchronous pipelines with the same
prefetch/overlap/spin switches as the TensorRT pipeline. The
replication tables use short rung names: \texttt{seq} is the sequential
pipeline, \texttt{prefetch} adds decode prefetch, \texttt{full
w1}/\texttt{w2}/\texttt{w3} add post-processing overlap with one, two
or three decode workers, and \texttt{spin} replaces the sleeping wait.
Pi~5 power is the rail sum of its power-management IC (PMIC), which
covers the system-on-chip (SoC) but not the 5\,V loads (the Hailo-8
itself, the fan, USB), so Pi energies are SoC-side lower bounds.

The Rubik Pi~3 (Ubuntu 24.04.4, kernel 6.8) is measured with two
stacks. The CPU stack is ONNX Runtime~1.30.0 (default thread
settings)~\cite{onnxruntime} on the same YOLO26n ONNX graph as the
Jetson float engine. The NPU stack runs an INT8 QNN context binary of
the same graph on the Hexagon HTP, with the standard
$(1{,}84{,}8400)$ head and host NMS at the val settings. The board's
QNN~2.46 predates the QAIRT~2.49.40-built context, so the context runs
through the bundled runtime of the \texttt{qai-appbuilder}~2.50.40
wheel. This userspace is not robust under concurrent calls, so all NPU
inference is serialized on one thread and host stages overlap around
it. The board has no trustworthy power sensor; its sampler logs only
the three CPU frequency domains and the SoC thermal zone, at an
effective 13--18\,Hz rather than 20\,Hz. The Pi sampler reads the
PMIC's SoC-side rails at an effective 10\,Hz. Energy integrals are
trapezoids over the actual sample timestamps. The HTP stack is
bit-deterministic (hash-verified, as on the Jetson and Pi). The
multi-threaded CPU stack emits rows in nondeterministic order, so its
equivalence is checked on sorted prediction rows and by COCOeval on
each rung's saved predictions.

\subsection{Models and engines}

The main model is the pretrained YOLO26n
checkpoint~\cite{jocher2026yolo26} (2.57\,M parameters); generality is
tested with YOLO26s (10.0\,M) and YOLOv10n
(2.78\,M)~\cite{wang2024yolov10}. All are exported to ONNX with
\texttt{nms=False}, which yields a $(1,300,6)$ tensor of boxes, scores
and class indices with no NMS. Ultralytics 8.4.172 loads the YOLO26
head with the one-to-many branch by default, so \texttt{nms=False} was
verified to select the one-to-one (end-to-end) branch in both export
and prediction. Every engine is a
static batch-1 FP16 TensorRT engine built on the device with builder
optimization level~3~\cite{nvidia2024tensorrt}, except the
compression-screening engines of \cref{sec:model}, built at level~0
(level~3 reference rows are marked in \cref{tab:compression}). Two
engines are built from each graph:
\begin{itemize}
\item a \emph{float} engine that takes a normalized
  $1{\times}3{\times}S{\times}S$ FP32 RGB tensor, and
\item a \emph{uint8} engine whose graph begins with four ONNX operators
  (Cast, a BGR$\to$RGB Gather, Transpose, division by 255) and takes the
  letterboxed $S{\times}S{\times}3$ uint8 BGR image directly; its ONNX
  Runtime outputs are identical to the float graph's (maximum absolute
  difference 0 on test images).
\end{itemize}
The Ultralytics baseline uses an engine exported on the device by
Ultralytics from the same checkpoint (\texttt{half=True, nms=False})
with the TensorRT builder defaults (optimization level 3, as above). It
runs through \texttt{YOLO.predict} with confidence threshold 0.001 and
at most 300 detections, as the custom pipelines do.

\subsection{Pipelines}
\label{sec:pipelines}

All custom pipelines share one implementation and differ only in the
switches below, which define the cumulative ladder of
\cref{sec:ladder}. Each goes from a JPEG file to COCO-format rows:
OpenCV decode, letterbox to $S{\times}S$ with value-114 padding,
copy into a pinned host buffer, TensorRT execution, output copy, and
vectorized conversion of rows with confidence $\geq 0.001$ to
original-image coordinates.
\begin{itemize}
\item \emph{Input type}: float input with host normalization, or uint8
  input with in-graph preprocessing.
\item \emph{Synchronization}: \texttt{cudaStreamSynchronize} on the
  inference stream, or \texttt{cudaEventSynchronize} on an event
  recorded after the output copy. \Cref{sec:sched} also
  sets the CUDA device scheduling flag (\texttt{auto}, \texttt{spin} or
  \texttt{blocking-sync}) before context creation.
\item \emph{Prefetch}: one or two worker threads decode and letterbox
  upcoming images while the current image is processed.
\item \emph{Overlap}: double-buffered I/O, so rows of image $i{-}1$
  are built while image $i$ runs on the GPU. The \emph{arrival-aware}
  variant (\cref{sec:stream}) finishes the in-flight image immediately
  if the next frame is not yet available.
\item \emph{Garbage collection}: Python's cyclic collector runs by
  default; one diagnostic variant (\cref{sec:gc}) disables it for the
  timed window.
\end{itemize}

\subsection{Workloads}

\emph{Backlog}: all images are available at once and processed as fast
as possible, giving throughput and energy per image. \emph{Streaming}:
frames become available at a fixed 15, 30 or 60\,fps over the first
1,000 val2017 images. Frame latency (\cref{tab:measure}) includes any
queueing when the pipeline cannot keep up. The median (p50) and the p95 are reported; with
1,000-frame streams the p99 estimate rests on ten samples, and p95 is
the highest quantile considered stable at this size. In both workloads
the first 20 images are a warm-up and are excluded. In streaming runs
the arrival schedule is re-anchored at the end of the warm-up, so the
timed window starts with an empty queue. JPEG files are read once
before each phase, so all decodes come from the page cache.

\subsection{Measurement boundaries}

\begin{table}[t]
\centering
\caption{Measurement boundaries used throughout. The terms differ in the
host work and queueing they span and are never equated.}
\label{tab:measure}
\small
\footnotesize
\begin{tabular}{@{}ll@{}}
\toprule
Term & Span (start $\to$ end) \\
\midrule
GPU compute & device-side execution of one inference; \\
\quad (\texttt{trtexec}) & CUDA events, no host work \\
Enqueue-to-done & H2D enqueue $\to$ completion of its \\
\quad (``GPU ms'' in tables) & D2H; includes GPU queueing \\
Accelerator call & submit $\to$ completion of one NPU \\
\quad (``infer'', NPU tables) & call; incl.\ queue wait in overlaps \\
Frame latency & frame availability $\to$ finished \\
\quad (p50/p95, ``latency'') & COCO rows; backlog: decode start \\
Completion period & $1/{}$throughput over timed window \\
Energy per image & $\int P\,dt$ over window $\div$ timed images \\
\bottomrule
\end{tabular}
\end{table}

\Cref{tab:measure} fixes the boundaries of the recurring time
quantities. The model's $G$
(\cref{sec:sysmodel}) is the device-side kernel cost, measured as
\emph{GPU compute} in the engine sweep and as \emph{enqueue-to-done}
inside the pipelines. On this unified-memory board the two differ only
by the launch and copy path: 0.1--0.4\,ms at the top clock (4.25\,ms
device-side versus 4.33--4.65\,ms enqueue-to-done for the same locked
engine). At low clocks the launch path itself stretches
(\cref{sec:sched}). The $M$ term of \cref{sec:sysmodel} is host work
outside $G$ on the pipeline's critical path and never double-counts
waiting already inside an enqueue-to-done span.

\subsection{Power, clock and thermal measurement}

A separate process samples \texttt{VDD\_IN} voltage and current, the GPU
\texttt{devfreq} clock, the highest CPU core clock and the junction
temperature (\texttt{tj-thermal}) at 20\,Hz. Energy is the trapezoidal
integral of power over each run's wall-clock timed window.
\texttt{VDD\_IN} is the module's main input rail (INA3221 monitor);
carrier-board loads (fan, USB peripherals, display) sit on other rails
and are excluded. Idle power was measured twice per clock policy over
60\,s windows (\cref{sec:ladder}). Unless stated otherwise, reported GPU
clocks are arithmetic means of the samples; \cref{eq:gf} uses the
harmonic mean. To check sampler overhead, the sequential and full
pipelines were each run twice with and without the power sampler on the
first 2,000 val2017 images, with a lighter GPU load and clock logger
running in both cases. Mean throughput differed by less than 0.2\%
(40.2 versus 40.2 and 209.7 versus 209.4 images/s), well within
run-to-run variation.

\subsection{Accuracy protocol and data hygiene}

Accuracy is official \texttt{pycocotools} COCOeval (bbox, maxDets 100)
on all 5,000 COCO val2017 images~\cite{lin2014coco}, reported as AP
(IoU 0.50:0.95), AP$_{50}$ (IoU 0.50) where noted, and AP$_S$ (small
objects). Every prediction file is hashed. Pipelines that differ only
in scheduling produce byte-identical predictions in every repeat and
under both clock policies, so their accuracy is exactly constant. Every
development decision (compression screening, router fitting, router
threshold) used two disjoint 1,024-image subsets of COCO train2017: a
development split and an independent, hash-verified verification
split. The pretrained checkpoint was itself trained on train2017, so
these splits are not unseen by the base model; they are held out only
from this paper's fitting and selection, their absolute AP is
optimistic, and they serve only for relative comparisons between
variants. val2017 was used only for the measurements in this paper.

\subsection{Use of AI assistance in the measurement process}
\label{sec:aiassist}

An AI coding agent (Factory Droid) ran the campaign scripts on the
boards, transferred the traces and drafted the aggregation and
plotting code. All experiments, protocols and analysis decisions were
made by the authors, and every script was reviewed before it ran. All
tables and figures are regenerated from the archived raw traces by the
published scripts (\cref{sec:data}).

\subsection{Repetitions and statistics}

Each backlog configuration on YOLO26n is repeated three times per clock
policy. Resolution, generality, streaming and scheduling-flag
experiments are repeated twice, and locked-clock generality runs once.
Follow-up experiments from later sessions used the same software and
protocol. Driver-reported GPU load was logged in one
run per pipeline. The utilization-clamping, GIL switch-interval,
minimum-frequency CPU clamp and float-engine zero-gap \texttt{trtexec}
cells have two runs each, and the fixed-frequency sweep has three 15\,s
runs per devfreq step (\cref{sec:diagnosis}). These runs also logged
GPU load, GPU clock and the clock of CPUs 0--3 at 20\,Hz.
Configuration order is rotated between repeats, with a 10\,s pause
between runs, and no other GPU work runs during a measurement. Means and full ranges (minimum to maximum) are reported
rather than confidence intervals, because two or three repeats are too
few for a meaningful interval. Run-to-run ranges are small (median 0.8\%
of the mean for backlog throughput). The held-out router test uses a
paired image bootstrap.
